% !TEX root = ./main.tex
\documentclass[a4paper]{article}
% Course language: use [english] or [german].
\usepackage[german]{ajnotes}
\title{Financial Accounting — Woche 3}
\begin{document}
    \maketitle
    % Attach the given problem sheet: drop it in this folder as sheet.pdf and
    % un-comment the next line (pdfpages is already loaded by the preamble).
    % \includepdf[pages=-]{sheet.pdf}
    \exercise{13.1} Start Up Bilanz
    \textbf{Kasse Konto (\textit{aktiv})} gegeben
        \begin{itemize}
            \item Soll (+)
                \subitem AB (3000) 
                \subitem (2) 8000
            \item Haben (-)
            \item (1) 3000
            \item 
        \end{itemize}
    \exercise{13.2} Unternehmen B
        Gegen Rechnung = Erhalt einer Ware aber noch keine Zahlung
        \textbf{Konto Kreditoren} passiv Konto gegeben
        \begin{itemize}
            \item Soll (-)
                \subitem  
            \item Haben (+)
                \subitem  
        \end{itemize}

        Debitoren schulden Kunden Geld
        Kreditoren schuldet die eigene Organisation Geld

    \exercise{13.4} Gliederung der Bilanz \linebreak
    Datum: 31.12.20.1 \linebreak
    \textbf{Konten}
    \begin{table}[]
        \begin{tabular}{llll}
            Kasse & Aktiv & UV & 10k \\
            Debitoren & Aktiv  & UV & 120k \\
            Immobilien & Aktiv  &  AV & 300k  \\
            Vorräte & Aktiv & UV  &  20k \\
            Darlehensschuld
        \end{tabular}
    \end{table}
    
\end{document}
